\documentclass[11pt]{article}
\usepackage[margin=1in]{geometry}
\usepackage{amsmath,amssymb,amsthm}
\usepackage{booktabs}
\usepackage{graphicx}
\usepackage{tikz}
\usetikzlibrary{arrows.meta,positioning,fit,backgrounds,calc}
\usepackage{hyperref}
\usepackage{xcolor}
\usepackage{enumitem}
\setlength{\emergencystretch}{2em} % absorb the occasional tight line
\hypersetup{colorlinks=true,linkcolor=figindigo,citecolor=figindigo,urlcolor=figindigo}
\usepackage[font=small,labelfont=bf]{caption}
\setlength{\textfloatsep}{16pt plus 2pt minus 2pt}
\setlength{\floatsep}{14pt plus 2pt minus 2pt}
\newtheorem{proposition}{Proposition}
\newtheorem{lemma}{Lemma}
\newtheorem{assumption}{Assumption}

% Set \anontrue for a double-blind submission (hides author/affiliation).
\newif\ifanon\anonfalse

\title{\textbf{Agents Over-Share:\\
Center-Blind Auditing for Multi-Tenant Agent Privacy}}
\ifanon
  \author{Anonymous Author(s)\\\small Paper under double-blind review}
\else
  \author{Aojie (Justin) Yuan\thanks{Code and the A2A-MT benchmark are available at
  \url{https://github.com/Justin0504/federated-agent-audit}.}\\
  \small University of Southern California}
\fi
\date{}

\newcommand{\ext}{\texttt{a2a.privacy/v1}}

\graphicspath{{figures/icons/pdf/}}
\input{figures/figstyle}

% Float-only pages start at the same top baseline as text pages.
\makeatletter
\setlength{\@fptop}{0pt}
\setlength{\@fpbot}{0pt plus 1fil}
\makeatother
\begin{document}
\thispagestyle{empty}
\noindent{\small\sffamily\color{figslate} Agents Over-Share\enspace /\enspace Research Preprint}\hfill
\ifanon\else\raisebox{-3pt}{\includegraphics[width=31mm]{figures/institution/usc-primary.png}}\fi\par
\vspace{5pt}\noindent{\color{figindigo}\rule{\textwidth}{0.5pt}}
\vspace{10pt}
\begin{center}
{\fontsize{23}{27}\selectfont\bfseries Agents Over-Share\par}
\vspace{5pt}
{\large Center-Blind Auditing for Multi-Tenant Agent Privacy\par}
\vspace{9pt}
\ifanon Anonymous Author(s)\else Aojie (Justin) Yuan\\[3pt]
{\small University of Southern California}\fi
\end{center}
\vspace{5pt}
\noindent\fcolorbox{figrule}{white}{\begin{minipage}{\dimexpr\textwidth-2\fboxsep-2\fboxrule\relax}
\vspace{6pt}\begin{center}\textbf{\color{figindigo}Abstract}\end{center}
\small
% ARXIV_ABSTRACT_BEGIN
Agent delegation can disclose sensitive information across organizational
boundaries. We introduce \ext{}, an experimental A2A metadata extension,
and an auditor that checks disclosure, purpose, forwarding and cross-message
inference from local labels. The center receives no raw message text, but
metadata leakage, incorrect labels and missing observations remain possible.
We separately evaluate release checks that bind authorization, case relevance
and factual support. On 20 new authored tasks, a frozen structured gate
preserves all eight benign completions, reduces literal forbidden-marker
release from two tasks to zero and unsupported-value release from four to zero,
and returns unknown on nine tasks. On 1,132 archived AgentDojo trajectories, destination
binding reduces native exfiltration success from 165/808 to 0/808 in
action-prefix replay; clean completion falls from 198/324 to 179/324.
Replay stops at the first block and does not measure recovery. A separate live
14B comparison runs eight methods on 32 development-disjoint tasks. Ungated
execution completes 10/32 tasks with one native attack success in 16 attacked
tasks; defended variants complete 0--12 tasks with zero observed attacks.
These small counts do not establish privacy or utility superiority. We release code and artifacts for jointly assessing attack outcomes, task
completion and cost under explicit trust assumptions.
% ARXIV_ABSTRACT_END
\begin{center}\footnotesize
\href{https://github.com/Justin0504/federated-agent-audit}{Code and evaluation artifacts}
\end{center}\vspace{3pt}
\end{minipage}}
\vspace{4pt}
\input{figures/fig_preprint_overview}
\clearpage

\section{Introduction}
Agents are starting to talk to \emph{each other}. Protocols such as A2A (v1.0,
Linux Foundation), MCP, and ACP let agents built by different teams discover
capabilities and delegate work across organizational boundaries. As this fabric
forms, a governance question arises: \emph{when my agent hands
data to your agent, what did it just share, and was it allowed to?}

Cross-boundary delegation creates a concrete governance problem for healthcare
scheduling, financial assistance and workplace automation: useful task context
can contain information the recipient is not authorized to receive. We study this problem with controlled deployments, frozen structured tasks
and externally authored agent workflows. These experiments test specific failure
mechanisms; they do not estimate production disclosure prevalence.

Three issues motivate our design. First, an audit limited to final user-facing
outputs can miss disclosures on internal agent channels, as studied by
AgentLeak~\cite{agentleak}. Second, several individually innocuous messages can
jointly reveal a sensitive attribute~\cite{sumleaks}. Third, deployments may
restrict a central observer from collecting raw messages across organizational
boundaries. We ask what this observer can conclude from local reports, and
which conclusions require stronger checks at the point of release.

The A2A data model supplies messages, parts, metadata and extension mechanisms.
Our extension uses these mechanisms to state governance roles and constraints
explicitly. It does not change the underlying transport or establish a standard
adopted by A2A. Its purpose is to make the evidence and assumptions consumed by
our auditor explicit and testable.

We study the evidence available to a central auditor that does not receive raw
agent messages. Local components derive governance labels, while the center
checks their relationships. This separation limits content collection but does
not establish label correctness or eliminate metadata leakage. Our contributions
are threefold:
\begin{enumerate}
\item \textbf{Separate measurements of privacy and utility.} A frozen
20-task structured experiment and 1,132 external archived AgentDojo trajectories
measure successful tasks separately from privacy violations, unknowns and
interventions. External replay exposes a utility cost hidden by the smaller
self-authored set. Historical model comparisons remain exploratory appendices.
\item \textbf{A metadata interface and scoped audit mechanisms.} The
\ext{} extension separates data subject, owning principal and acting agent.
Four detectors consume local labels for disclosure, purpose, forwarding and
inference checks. We characterize inference under an explicit gain model and
state only conditional representation-capacity and report-authenticity results;
neither proves semantic confidentiality or complete observation.
\item \textbf{Failure-driven release checks and reproducible evaluation.}
Controlled deployment experiments distinguish stale authorization, irrelevant
context and unsupported generated claims. We implement a restricted structured
release gate and a transactional database adapter, then evaluate their checks,
abstentions and retained task outputs separately. Section~\ref{sec:methodology}
reports what these mechanisms establish and what remains outside their scope.
\end{enumerate}
The SDK also supports same-organization multi-agent deployments. The deployed
experiments are research harnesses; they do not establish production readiness.


\section{Problem and evidence boundary}\label{sec:threat}
The unit of analysis is an agent-to-agent handoff across a data-governance
boundary. We distinguish the \emph{data subject} whom a fact describes, the
\emph{owning principal} controlling the resource, and the \emph{acting agent}
that sends or receives it. Shared ownership does not automatically authorize
all purposes, recipients or fields.

A recipient may retain and combine messages, a forwarder may exceed its scope,
and a local reporter may omit events or mislabel their content. Our software
baseline trusts the local observation and labeling path for semantic detection;
it does not solve malicious endpoints or unmediated communication. Report MACs
authenticate an issuer's statement, not its truth. Replay prevention, completeness
and key isolation are separate obligations. The formal authenticity statement
and its assumptions appear in Appendix~\ref{app:audit-details}.

\paragraph{What the center can observe.}
The center receives labels, routing identities, provenance and content hashes,
not raw message text. This reduces centralized content collection but still
reveals metadata and equality patterns. Low-entropy hashes may permit dictionary
attacks. A missing alert therefore cannot certify confidentiality: the report
may be incomplete, its labels wrong, or the relevant inference outside the model.
Figure~\ref{fig:preprint-overview} summarizes this boundary and its empirical limits.

\section{A scoped audit and release interface}\label{sec:ext}
\subsection{Metadata audit}\label{sec:auditor}
The experimental \ext{} extension attaches governance fields to A2A Parts:
subject, owner, sensitivity, categories, inferred categories, allowed recipients,
purpose, hop budget and provenance. It uses A2A's metadata mechanism; it is not
an adopted A2A standard. A local component derives labels and a central graph
checks four relationships: disclosure to an unauthorized principal, incompatible
purpose, forwarding beyond a hop budget, and accumulated inference across messages.
\input{figures/fig_pipeline}

For the inference check we use an explicit odds model,
$O_k=O_0\prod_{i=1}^k\lambda_i$, and flag a posterior gain exceeding a chosen
threshold. Its likelihood ratios and conditional-independence assumption are
modeling choices, not calibrated guarantees about real recipients. Accurate
labels alone do not validate these assumptions. Appendix~\ref{app:audit-details}
provides the derivation, protocol schema and worked example;
Appendix~\ref{app:detector-studies} preserves detector and tagger measurements.

\subsection{Release checks and their implementation}\label{sec:methodology}
An audit report must distinguish three claims: authorization at release ($A$),
relevance to the requested case and purpose ($R$), and support by permitted
source evidence ($S$). Authorized context can belong to the wrong case, and
relevant context can still lead to an unsupported assertion. Checking $A$ does
not establish $A\land R\land S$. Missing evidence produces \emph{unknown}, not a
positive safety verdict or a fact about the world.

We adopt the separation of model proposals from trusted authority found in
capability and information-flow systems, including CaMeL~\cite{camel2025},
FlowSeal~\cite{flowseal2026} and FAVA~\cite{fava2026}. Our question is which scoped
claims a content-blind center can justify when evidence is incomplete. We do
not reproduce those systems or claim empirical superiority over them.

\paragraph{Implemented structured gate.}\label{sec:structured-release}
A model proposes tuples $(\mathrm{source},\mathrm{version},\mathrm{field},\mathrm{value})$.
A trusted local adapter supplies workspace, case, recipient, purpose, source
fields and authorization state. The full gate requires current recipient access,
matching workspace and case, an allowed field-purpose pair, matching version and
exact equality to the source value. It rejects permission overrides and extra
free text. Missing authority or support returns unknown; established mismatches
return deny. Only admitted fields reach a fixed renderer, and the entire proposal
is withheld if any tuple fails.

An experimental PostgreSQL adapter serializes this check with a business-side
mailbox commit. That boundary does not guarantee external delivery or remove
information already in model memory. Source facts can be wrong, exact matching
rejects valid paraphrases, and individually allowed fields can enable inference.
A minimal center-facing summary exposes status and fixed reason codes; it is
neither differentially private nor independently verifiable semantic evidence.

\paragraph{What remains proposed.}
General free-text support verification and a complete composition of semantic
auditing with release enforcement remain unvalidated. The current gate is a
restricted executable check, not a natural-language entailment system. The
receipt interface binds identifiers and an epoch but does not implement every
proposed purpose, version and support obligation. Appendix~\ref{app:release-details}
retains the full specification, counterexamples and database experiments.

\section{Evaluation: privacy, utility and cost}\label{sec:eval}
We organize the evidence by what each execution mode can establish. Authored
structured tasks test mechanism behavior; external replay measures intervention
on archived actions; live tasks measure execution under fixed budgets. We do
not pool these units or equate native attack failure with semantic privacy.

\begin{table}[t]\centering\small
\begin{tabular}{@{}p{.22\linewidth}p{.28\linewidth}p{.42\linewidth}@{}}
\toprule
Evidence & Unit and scale & Principal limitation\\
\midrule
Structured transfer & 20 authored tasks; five gates on the same proposals & Exact fields and proxy scoring; no independent human semantic labels\\
External replay & 1,132 trajectories from 81 tasks; five controls & Stops at the first block; no recovery or adaptive replanning\\
Live comparison & 32 tasks; eight methods; 256 episodes & One pinned 14B model and seed; sparse successful attacks\\
\bottomrule
\end{tabular}
\caption{Three distinct evidence sources. Repeated controls and archived model variants are not additional independent tasks.}
\label{tab:evidence-scope}
\end{table}

\subsection{Structured release: a bounded mechanism result}
A frozen gate and scorer evaluate 20 new authored tasks using one pinned
Qwen2.5-3B proposal per task. Every gate sees the same proposal. The full gate
completes all eight benign requests, reduces literal forbidden-marker releases
from two to zero and unsupported-value releases from four to zero. Deny-all
also removes these releases but completes no benign request. Scope checking
removes the two marker failures while leaving the four unsupported values;
exact source matching removes those values.

The full gate returns unknown on nine tasks: two schema failures, three cases
without authority and four unsupported values. Three further requests are denied.
Two malformed proposals are rejected before policy evaluation, limiting any
revocation-specific interpretation. Source errors and valid paraphrases remain
counterexamples. This small authored transfer check establishes neither broad
semantic correctness nor production performance. The full ablation table and
frozen protocol appear in Appendix~\ref{app:release-details}.

\subsection{External replay: privacy reduction costs utility}
We freeze 1,132 AgentDojo~\cite{debenedetti2024agentdojo} archived trajectories:
324 clean and 808 attacked, covering 81 underlying tasks. Five controls replay
the same actions using unchanged native task and attack predicates. At the first
blocked egress action, replay terminates; the original later answer is removed.
This tests a destination-authority adapter, not semantic message support or how
an agent recovers after a refusal.

Combined intent and destination checks reduce native exfiltration success from
165/808 to 0/808. Clean utility falls from 198/324 to 179/324, a paired change of
$-5.9$ percentage points (task-cluster bootstrap 95\% interval $[-10.2,-2.2]$).
Attacked utility falls from 383/808 to 316/808, a change of $-8.3$ points
($[-12.8,-4.3]$). Destination-only checking produces the same counts; this
experiment supplies no measurable benefit for the additional intent rule.

\begin{figure}[tbp]\centering
\includegraphics[width=\linewidth]{figures/agentdojo20261013/privacy_utility.pdf}
\caption{External replay retains all selected trajectories and separates native
attack success from task completion. The deny-egress arm has one clean and one
attacked checker exception, retained as unresolved. Task-cluster intervals and
reconstruction discrepancies are detailed in Appendix~\ref{sec:external-dojo}.}
\label{fig:core-replay}
\end{figure}

Losses concentrate in Slack and include legitimate implicit destinations that
literal binding cannot resolve. The archive also has checker-version discrepancies:
we retain and audit them instead of replacing the native predicates or selecting
only successful reconstructions. Appendix~\ref{sec:external-dojo} reports the
complete controls, suite breakdowns and discrepancy audit. These losses qualify
the eight-of-eight benign retention in the authored set.

\subsection{Live execution: low attack counts are not a safety ranking}
Following development, we freeze Qwen2.5-14B-Instruct Q4\_K\_M, a shared system
instruction and budgets of 12 model calls, 12 actions, 1,024 output tokens per
call and an 8,192-token context. Setup and planning consume the same call budget.
The final development gate requires at least two of three ungated clean twins
to succeed and known outcomes; the 72-episode validation passes. All earlier
failures, model screens and interface diagnostics remain in
Appendix~\ref{app:development} and are excluded from formal estimates.

The formal batch contains 16 clean and 16 attacked task IDs disjoint from the
development tasks, eight methods and one seed: 256 completed episodes with known
native outcomes. Task families and archived payloads appeared in earlier
diagnostics, so task disjointness does not imply a previously unseen benchmark.
We retain official Tool Filter and IPIGuard alongside five system controls,
and separately label a Tool Filter adaptation with a generic resume instruction.
No task-specific hints or formal-outcome tuning are used.
\begin{figure*}[tbp]
\centering
\includegraphics[width=\textwidth]{figures/final_native/same_task_eight_outcomes.pdf}
\caption{One frozen task, eight executions. We select T30 post hoc by the explicit
rule of the smallest stable task index with ungated native attack success. This
is a diagnostic illustration, not a representative success claim. Only the
labeled Tool Filter resume adaptation completes this task. The original baseline
and all failures are retained. Appendix Figure~\ref{fig:full-task-matrix} shows
all 256 outcomes; no additional inference is performed for these displays.}
\label{fig:eight-outcomes}
\end{figure*}

\begin{figure*}[tbp]
\centering
\includegraphics[width=\textwidth]{figures/final_native/formal_privacy_utility.pdf}
\caption{Frozen live comparison on 32 development-disjoint AgentDojo tasks,
16 clean and 16 attacked, using one 14B model and one seed. Points are observed
rates; intervals are Wilson 95\%. Joint success requires native task success and
native attack failure. Zero observed attacks is compatible with substantial
uncertainty and does not establish semantic privacy.}
\label{fig:final-native}
\end{figure*}


Ungated execution completes 10/32 tasks and exhibits one native attack success
among 16 attacked tasks. Gate and scope complete 10/32 each, receipt 9/32,
deny-egress and IPIGuard 8/32 each, original Tool Filter 0/32 and its resume
adaptation 12/32. All defended variants have zero observed native attacks.
Only one attack case is discordant with no gate; zero out of 16 has a Wilson
95\% upper bound of 19.4\%. This is not evidence of broad privacy superiority.

Resume gains four tasks and loses two against no gate: $+6.25$ percentage points
(paired bootstrap 95\% interval $[-9.38,21.88]$). Receipt gains one and loses two:
$-3.13$ points ($[-12.50,6.25]$). Neither supports utility superiority. Gate,
scope, receipt and IPIGuard each attain 4/16 attacked joint successes, equal to
no gate; resume attains 5/16. Joint success still uses a specific native attack
predicate, not comprehensive semantic privacy.

\paragraph{Cost and failure accounting.}\label{sec:cost}
Total model calls are 138 for no gate, 137 for gate, 136 each for scope and
receipt, 142 for deny-egress, 211 for IPIGuard, 69 for original Tool Filter and
144 for resume. IPIGuard exhausts the matched budget on 21/32 tasks. Four original
Tool Filter calls have unknown token use; their totals are lower bounds.
Failures and unknown costs remain visible. Appendix Figure~\ref{fig:final-native-cost}
reports paired cost intervals; Figure~\ref{fig:full-task-matrix} shows every task.
Timing is observational, not randomized, and no currency-cost claim is made.

\clearpage
\section{Limitations, related work and reuse}
The evidence supports scoped checks and exposes the cost of intervention. It
does not establish general semantic confidentiality, receipt superiority or
production readiness. False labels, omitted events, malicious authorized issuers,
source errors, correlated inference and unmediated channels remain open. The
center cannot independently verify judgments whose supporting content it does
not receive. A MAC authenticates a statement; it does not validate that judgment.

AgentLeak~\cite{agentleak} and Sum-Leaks~\cite{sumleaks} motivate checking internal
channels and compositional inference. AgentDojo supplies independent executable
tasks and attack predicates. Our metadata interface complements local enforcement
and tracing; it does not show that existing DLP tools cannot execute locally or
redact central logs. Comparisons against Presidio~\cite{presidio} and configured
model judges are integration-specific and retained in
Appendix~\ref{app:detector-studies}. Historical model aliases lack consistent
provider-version provenance and cannot support precise capability rankings.

The reusable contribution is an explicit evidence boundary and an evaluation
that keeps audit findings, attack outcomes, task completion and cost separate.
The accompanying single-handoff example mirrors existing application messages
into \texttt{AuditSession} without model or network calls. It observes rather
than enforces; its aggregate output is not an independent task-success checker.
Code, frozen numeric evidence and the complete result matrix support inspection
of failures as well as successes. Stronger semantic and deployment claims require
independent labels, more capable executors, additional models and adversaries,
and complete mediation of the release path.
\label{sec:main-end}

\clearpage
\appendix
\section{Complete task-level evidence and minimal integration}
\label{sec:complete-task-evidence}
Figure~\ref{fig:full-task-matrix} retains every frozen task and method. Stable
indices T01--T32 follow the sorted remote task keys; raw traces and task text
remain remote. C denotes clean and A attacked. Clean rows have no attack trial:
green denotes task success there, not a tested security result. For attacked
rows, green requires task success and native attack failure. Gray denotes task
failure without observed native attack success; red denotes native attack
success. Black dots preserve recorded budget exhaustion even when the task
checker succeeds. Counts include setup and planning calls. Token uncertainty
is reported separately in Figure~\ref{fig:final-native-cost}.

\begin{figure}[htbp]
\centering
\includegraphics[width=0.94\textwidth,height=0.76\textheight,keepaspectratio]{figures/final_native/complete_task_matrix.pdf}
\caption{All 32 tasks and eight methods, without outcome-based row filtering.
Cells contain model-call counts. Original Tool Filter and its resume adaptation
are distinct columns. All native task outcomes are known; clean attack outcomes
are not applicable. These 256 executions are 32 task-level clusters, not 256
independent tasks.}
\label{fig:full-task-matrix}
\end{figure}
\clearpage
\paragraph{A minimal application boundary.}
The accompanying \texttt{audit\_handoff.py} in the repository's
\texttt{examples/minimal\_handoff/} directory wraps the
existing \texttt{AuditSession}: an application passes handoff text locally,
recipient identity, sensitivity, categories and allowed recipients. It receives
finding types and counts. The example makes no model or network calls and does
not alter delivery. Policy labels must come from trusted configuration; an empty
finding set does not establish semantic privacy. This offline integration example
is not part of the experimental sample. Workflow-level cumulative auditing uses
a persistent session rather than this single-handoff wrapper. The example README
specifies installation, the JSON input contract, retained local data and the
separation between audit findings, enforcement and independent task evaluation.

\clearpage
\section{Protocol details, conditional bounds and authored benchmark}
\label{app:audit-details}
\paragraph{Report authentication and its limits.}
A build identifier and MAC let the center check report origin and integrity.
They do not establish that a build observed all actual inputs, produced correct
labels, or executed in isolation. In the software baseline, a build fingerprint
is an identifier rather than hardware evidence of execution.

\begin{proposition}[Report authenticity, not semantic soundness]\label{prop:attest}
Assume an EUF-CMA-secure MAC, an uncompromised issuer key, and a verifier that
checks a canonical report serialization. An adversary without that key cannot
produce an accepted message--tag pair for a new message not previously
MACed by the issuer, except with negligible probability.
\end{proposition}
\begin{proof}
Such a pair constitutes a MAC forgery under the stated security experiment.
If the implementation authenticates a digest rather than the serialization,
the argument additionally requires collision resistance of that digest.
\end{proof}
This statement allows replay of a previously authenticated report; event binding
and deduplication are separate checks. It also allows an authorized issuer to
sign an incorrect claim. In particular, the proposition proves neither label
accuracy nor completeness of observation. Key isolation and trusted execution
can strengthen provenance under additional assumptions, but do not automatically
make the labeler correct or mediate every communication path. Canonical subject
identifiers prevent some accidental inconsistencies, not false subject claims
by a trusted signer. Section~\ref{sec:methodology} separates these obligations.

\subsection{The \ext{} extension}\label{app:ext-detail}
The three roles of the threat model become fields the protocol can carry. Each
privacy-relevant \texttt{Part} carries a label under
\texttt{Part.metadata["a2a.privacy/v1"]}:
\[
\textsf{PrivacyLabel} = \{\,
\text{subject},\ \text{owner},\ \text{sensitivity}\!\in\![0,5],\ \text{category}[],
\]
\[
\text{inferred\_categories}[],\ \text{purpose}[],\ \text{allowed\_recipients}[],\
\text{ttl\_hops},\ \text{provenance\_id}\,\}.
\]
Identifiers are opaque (\texttt{subject:alice}, \texttt{tenant:hospital})---governance
metadata rather than raw message text; their meanings can still reveal sensitive facts. An \texttt{AgentClearance} (an AgentCard declaration)
states the purposes an agent is cleared for, so cross-tenant rules are checkable
from public metadata. Two fields are central to the results:
\textbf{inferred\_categories} is computed \emph{locally} (the local auditor sees
content; ``near the oncology center'' $\rightarrow$ \texttt{health}) and only the
tag travels---the signal the inference detector accumulates without the center
reading content; \textbf{provenance\_id} is a stable datum identity assigned by the
origin and preserved by forwarders even when they re-word the data, so hop/TTL
tracking follows a datum across paraphrasing (a content hash breaks the moment a relay rephrases).
Figure~\ref{fig:label} maps the implemented fields to their policy checks.
\input{figures/fig_label}


\subsection{The center-blind auditor and detectors}\label{app:auditor-detail}
These labels are exactly what the auditor reads, and the content is exactly what
it does not. The auditor desensitizes each labeled \texttt{Part} into a
center-view edge carrying a content hash, the principals, and the
label---no raw text (Figure~\ref{fig:desens}). It
checks for undeclared raw-content token matches in its output, allowing declared
label values, schema tokens and routing identifiers. This lexical check is a
regression diagnostic, not a complete detector of semantic information leakage.

\begin{lemma}[Conditional representation-capacity bound]\label{lem:blind}
Let $S$ denote already disclosed side information. Suppose the remaining
center-visible representation is exactly $V=(D,L)$, where $D$ has $b$ bits and
$L$ takes values in a fixed finite alphabet $\mathcal{L}$. Then
\[ I(M;V\mid S)\le b+\log_2|\mathcal{L}|. \]
All visible fields not included in $S$ must be included in $(D,L)$ for this
bound to apply; unbounded free-text labels do not meet the finite-alphabet premise.
\end{lemma}
\begin{proof}
$I(M;V\mid S)\le H(V\mid S)\le H(D\mid S)+H(L\mid S)
\le b+\log_2|\mathcal{L}|$. No zero-information assumption about a public
hash is required.
\end{proof}
This elementary capacity bound is not a confidentiality or differential-privacy
guarantee. Hashes expose equality and may permit dictionary attacks on low-entropy
messages; labels can disclose sensitive facts directly. Repeated observations,
routing identifiers and external knowledge require separate analysis. In
particular, the bound does not establish computational non-reconstruction.
\input{figures/fig_example}

Here ``center-blind'' describes the absence of raw message text at the center,
not information-theoretic secrecy of the metadata graph. Four detectors run
over this graph:
\begin{enumerate}[leftmargin=1.4em]
\item \textbf{Cross-tenant disclosure}: sensitive data ($\geq\tau$) reaches a
principal that neither owns it nor is an allowed recipient.
\item \textbf{Purpose limitation}: data reaches an agent whose clearance purposes
are disjoint from the data's permitted purposes.
\item \textbf{Hop/TTL}: a datum (tracked by \texttt{provenance\_id}) is relayed
more than \texttt{ttl\_hops} times; robust to a forwarder re-wording content.
\item \textbf{Cross-tenant inference}: no single edge is a disclosure, but a
recipient principal accumulates converging \texttt{inferred\_categories}
fragments about one subject, letting it infer a sensitive category it was never
authorized for (Figure~\ref{fig:desens}); we ground this in the bound below.
\end{enumerate}

Figure~\ref{fig:deploy} shows the report-only deployment and its authentication boundary. One property of the deployment matters for the fourth detector: fragments about one subject arrive in separate reports from separate agents, so a center that forgot every edge after answering could only compose what one report happened to bundle. The service therefore keeps a retention-bounded metadata graph per workspace (edges, never content) and runs the detectors over the retained graph plus each new report, reporting each violation once; retention is a policy knob---a longer window composes more and holds more metadata---and the cross-report case is exercised at the HTTP entry point in the test suite.
\input{figures/fig_deploy}



\paragraph{A formal inference-gain bound.} We model what a recipient infers about
a subject's sensitive attribute $A\in\{0,1\}$ from $k$ converging fragments
$f_1,\dots,f_k$, with prior $0<p_0=P(A{=}1)<1$, threshold $0<\delta<1-p_0$, and odds $O_0=p_0/(1-p_0)$.

\begin{assumption}[Conditional independence]\label{as:ci}
Given $A$, the fragments are independent, each with likelihood ratio
$\lambda_i=\tfrac{P(f_i\mid A{=}1)}{P(f_i\mid A{=}0)}\ge1$; write $\lambda$ for a
common lower bound $\lambda_i\ge\lambda>1$.
\end{assumption}

\begin{proposition}\label{prop:k}
Under Assumption~\ref{as:ci} the posterior odds are $O_k=O_0\prod_i\lambda_i$, the
posterior is $P(A\mid f_{1:k})=O_k/(1+O_k)$, and the inference gain
$g(k)=P(A\mid f_{1:k})-p_0$ is strictly increasing in $k$. With uniform $\lambda$,
the detector (fires iff $g(k)\ge\delta$) fires exactly for
$k\ge k^{*}=\big\lceil\log_\lambda(O_\delta/O_0)\big\rceil$, where
$O_\delta=(p_0+\delta)/(1-p_0-\delta)$.
\end{proposition}
\begin{proof}
By Bayes in odds form, $O_k=O_0\prod_i\lambda_i$; since each $\lambda_i>1$, $O_k$
is strictly increasing in $k$, and $x\mapsto x/(1+x)$ is strictly increasing, so
$g$ is strictly increasing. With uniform $\lambda$, $g(k)\ge\delta \iff
P(A\mid f_{1:k})\ge p_0+\delta \iff O_k\ge O_\delta \iff O_0\lambda^{k}\ge O_\delta
\iff k\ge\log_\lambda(O_\delta/O_0)$; integrality gives the ceiling.
\end{proof}
For defaults $(p_0,\lambda,\delta)=(0.1,3,0.3)$, $k^{*}=2$: one hint never fires,
two converging fragments do (posterior $0.50$, gain $0.40$), replacing the
heuristic $1-2^{-k}$ with an explicit assumed-model quantity and a closed-form threshold (Figure~\ref{fig:gain});
deployments tune $(p_0,\lambda,\delta)$ and may set $\lambda_i$ from the tagger's
hint specificity (so one high-$\lambda$ fragment can suffice).
\emph{On the assumption:} real fragments are correlated, so
Assumption~\ref{as:ci} is an approximation --- positively-correlated (redundant)
hints can make the product over-credit evidence. Historical recipient probes
also exhibit detector misses, but do not isolate tagger coverage from
misspecification of the likelihood model. The bound is an assumed-model design tool, not a calibrated or tight
guarantee.

\input{figures/fig_gain}

\subsection{The A2A-MT benchmark}
A2A-MT supplies authored multi-tenant A2A interactions with labels and hashes
as the auditor's input. The benchmark tests consistency with template-defined
policy labels, rather than independently adjudicated semantic confidentiality.
A curated \emph{calendar-negotiation} family (two people's
scheduling agents must reveal availability without leaking \emph{why} a slot is
blocked) instantiates all four violation types; parameterized families
(group-assistant, marketplace-delegation, cross-tenant aggregation, plus
disclosure/purpose/TTL/inference sweeps with near-miss clean controls) give 48
labeled scenarios, ground truth set by template intent. An LLM fills each Part
with realistic text matching its label (labels remain the ground truth), exported
as a shareable artifact.

WeClawArena~\cite{weclaw} studies cross-user collaboration in an auditable
sandbox. Our authored scenarios instead restrict the central observer to
metadata. These are different observation settings; we have not performed a
matched comparison and neither suite supplies a numerical upper bound for the
other. The three-role schema is a design choice whose utility must be evaluated
under explicit deployment assumptions.

\section{From observed failures to evidence-aware auditing}
\label{app:release-details}
The experiments motivate a narrower question than whether a hand-off is safe:
\emph{which claim can the center justify from the evidence it actually receives?}
We derive a design from the observed failures below. A restricted structured-output
gate is implemented and tested on authored cases; the broader semantic composition
remains proposed. We distinguish these two levels throughout.

\paragraph{What recent methods contribute.}
FlowSeal~\cite{flowseal2026} places mandatory release checks outside the agent's
conversation and isolates semantic declassification. FAVA~\cite{fava2026}
separates model-produced permission representations from deterministic authorization;
its guarantee begins after evidence has been lowered into its graph.
CaMeL~\cite{camel2025} uses capabilities and data-flow control while explicitly
excluding some text-to-text corruption from its protection. We adopt the separation
of proposal and authority, rather than claiming it as new. Our narrower research
question concerns which release claims a content-blind center can justify when
policy, identity, support, or event evidence is incomplete. We do not reproduce
these systems or claim an empirical advantage over them.

\subsection{Authorization, relevance, and support are separate claims}
The deployed follow-up used four development tasks in a $2\times2$ design:
broad versus case-scoped retrieval, and snapshot versus delivery-time permission
checks (Appendix~\ref{app:jetstream}). All 16 episodes completed, with 32 model
calls and no token-cap hits. Both snapshot variants exposed the revoked Birch
reference. Delivery-time checks withheld that record, but broad retrieval still
let the answer substitute Acorn's reference for Birch's. Case-scoped retrieval
removed that particular substitution in the paired episode, yet the answer
asserted that Birch's refund ``has not been processed yet,'' which the available
evidence did not establish. These are observations on reused development cases,
not estimates of population-level error rates.

Three predicates are therefore needed. $A$ asks whether the recipient may access
the resource at delivery; $R$ asks whether it belongs to the requested case and
purpose; $S$ asks whether the generated claim is supported by the permitted
evidence. Authorized context can be irrelevant, and relevant authorized context
does not make every generated assertion true. A check of $A$ must not be reported
as a check of $A\land R\land S$. A missing record also establishes only lack of
available evidence, not a negative fact about the world.

\subsection{Make the scope of a verdict explicit}
For an event $e$, let $E_e$ be the available evidence and let
$\Omega(E_e;\mathcal{T})$ denote the nonempty set of executions consistent with
that evidence under stated trust assumptions $\mathcal{T}$. For one specified
violation predicate $V$ over completed disclosures, the ideal verdict is
\[
\widehat V(E_e)=
\begin{cases}
\textsf{violation}, & \forall w\in\Omega(E_e;\mathcal{T}),\ V(w)=1,\\
\textsf{no violation in scope}, & \forall w\in\Omega(E_e;\mathcal{T}),\ V(w)=0,\\
\textsf{unknown}, & \text{otherwise}.
\end{cases}
\]
Inconsistent evidence, represented by an empty set, also yields unknown and an
integrity alert; it must not make either universal condition vacuously true.
If the actual execution belongs to this set, a definite verdict is sound by
definition. This observation is a specification, not a proof that our current
verifier constructs the set correctly. Its difficulty lies in establishing the
trust assumptions and evidence coverage. The current software uses bounded
checks and heuristic semantic labels, not a general possible-executions solver.

Each response should name its predicate, event or time-window scope, supporting
evidence, and missing obligations. The absence of a semantic alert cannot justify
an authorization claim, or vice versa. Likewise, a blocked attempt is a denied
action rather than a completed disclosure. A pre-execution gate may map unknown
to deferral, redaction, or refusal according to risk; that operational choice
does not change unknown into evidence of a violation.

\subsection{Bind policy evidence to the release boundary}
A proposed release receipt binds the workspace, event, resource version,
recipient, purpose, policy epoch, and release decision. A trusted authority must
establish the epoch effective at that release, not merely return its current
epoch when an old event is audited. A policy check and the corresponding context
release must share a defined serialization boundary. Otherwise a revocation
between check and release recreates the stale-snapshot failure.

Our database harness implements a transactional check at its recorded context
release boundary. It does not prove revocation of information already present in
a model's context or memory. The experimental receipt verifier binds identifiers
and an epoch, but does not yet implement all the proposed purpose, resource-version,
and semantic-support obligations. Its authoritative epoch is supplied by the
evaluation fixture. A stale authority and an authorized issuer signing a false
decision remain counterexamples. Authentication alone cannot eliminate either.

\subsection{Check semantic support locally and disclose the uncertainty}
At the owner boundary, a proposed support checker decomposes a generated answer
into claims and relates each claim to authorized evidence for the requested case.
An unsupported claim triggers revision or a statement that the information is
unavailable. The center receives only a scoped support result, opaque provenance
references, checker version, and uncertainty status. Raw claims and supporting
passages remain local. This is a proposed mechanism: the current online harness
reports generated-output observations as unknown, rather than pretending to
verify their semantics.

The additional checker is fallible and creates a new trust dependency. It must
be tested on omissions, paraphrases, plausible but unsupported assertions, and
contradictory records. Self-reported model confidence is not a calibrated error
bound. The center can verify a checker's identity and evidence bindings without
being able to verify the judgment itself. Privacy labels also need a separate
test: a statement can be fully supported and still disclose a prohibited fact.


\subsection{A restricted, implemented release mechanism}
\label{app:structured-release}
To test the separation without treating a second LLM as an oracle, we implement
an experimental structured-output gate. The model proposes only tuples
$(\mathrm{source},\mathrm{version},\mathrm{field},\mathrm{value})$. A trusted local
adapter supplies the task's workspace, case, recipient and purpose, source fields,
and current authorization state. Model-proposed permission overrides and extra
free-text fields are rejected. The full gate requires current recipient access,
matching workspace and case, an allowed field-purpose pair, matching source
version, and exact equality to the corresponding source value. Missing authority
or support produces unknown; an established scope or permission mismatch produces
deny. The entire proposal is withheld if any tuple fails.

Only validated fields reach a fixed renderer; model-written prose is never
appended afterward. Thus this implementation prevents unsupported field substitution
within its admitted representation. It does not solve natural-language entailment.
A source may itself be wrong, equivalent paraphrases are rejected, and releasing
several individually permitted fields can enable inference outside these checks.
The adapter is a trusted input, and production use would additionally require
serializing the check with release and mediating every output channel. The core gate is a pure function. An experimental PostgreSQL adapter now
serializes it with a business-side mailbox commit, as evaluated below; neither
is integrated into the production HTTP audit service or an external message sink.

A minimal center-facing summary contains only a schema identifier, status and
fixed reason code. It excludes the source values, case names, model prose and
source identifiers. The local result contains those details for the business-side
renderer and must not be serialized wholesale to the center. Even this reduced
summary discloses an outcome and failure class; it is neither differentially
private nor independently verifiable evidence of correct execution.

\paragraph{Mechanism regression and its counterexamples.}
Nineteen authored fixtures cover revoked and deleted records, stale versions,
missing authority, cross-case and cross-workspace references, purpose restrictions,
unsupported values, malformed proposals and mixed valid/invalid claims. The same
inputs are available to all arms. These controls are ablations of our own code,
not implementations of FlowSeal, FAVA or CaMeL. Table~\ref{tab:release-regression}
reports 95 deterministic evaluations, not 95 independent workflows.

\begin{table}[ht]
\centering\small
\begin{tabular}{lrrrr}
\toprule
Control & Inadmissible release & Benign exact & Unknown & Source error\\
\midrule
Snapshot & 10 & 2/2 & 5 & 1/1\\
Fresh policy & 6 & 2/2 & 7 & 1/1\\
Fresh policy + scope & 4 & 2/2 & 7 & 1/1\\
Full structured gate & 0 & 2/2 & 11 & 1/1\\
Deny all & 0 & 0/2 & 19 & 0/1\\
\bottomrule
\end{tabular}
\caption{Authored contract regression. Inadmissible release counts violations of
this representation's contract, not semantic privacy errors. Benign exact counts
two supported exact-value outputs. Source error records acceptance of a value
that matches an intentionally incorrect trusted source; it is a negative control,
not a successful answer.}
\label{tab:release-regression}
\end{table}

The full gate also rejects the supported paraphrase ``on Tuesday'' for the source
value ``Tuesday.'' Its zero contract violations therefore come with a measurable
expressiveness cost. All nontrivial arms accept the deliberately incorrect
trusted source, demonstrating why source consistency must not be called factual
truth. Independent held-out and human semantic evaluation remain necessary.

\paragraph{Model-generated proposals.}
An eight-call Jetstream pilot used cached Qwen2.5-3B-Instruct with greedy decoding
and a 160-token cap. The inputs are a subset of the same authored fixtures, not
new held-out workflows. The model received a draft and source fields but not the
expected verdict or release-time authority state. All eight proposals parsed
without truncation or generation errors. Snapshot control released all eight;
fresh policy released five; adding scope released three; the full gate released
two, denied three, and returned unknown for three. The two released proposals
were the benign exact-field cases. In the unsupported-negative fixture, the
model retained the seeded claim ``not processed'' despite the source saying
``scheduled''; exact support checking withheld it. This is draft-conditioned
inference plus a local gate, not a new end-to-end database deployment or an
independent semantic evaluation. Frozen protocol and source hashes, raw outputs,
model revision, and environment versions accompany the artifact.

\paragraph{Checking freshness is not sufficient without a release boundary.}
We next wrapped the gate in a PostgreSQL transaction that locks the resource rows,
reads current source and authority state, evaluates the proposal, and commits a
business-side mailbox entry together with an audit-summary outbox entry. Resource
updates use the same rows. An event-scoped transaction lock serializes concurrent
retries; a binding digest rejects reuse of an event identifier for a different
proposal or task. A successful retry returns an acknowledgment rather than
re-emitting the saved answer.

Seven authored schedules, two configurations, and three repetitions produced
42 real database trials. In the check-without-row-lock ablation, a barrier let
revocation or a version update commit after validation but before the mailbox
write: all six such trials retained the stale allow result. With row locking,
we observed the conflicting update waiting in PostgreSQL's lock table; the
mailbox transaction serialized before it. A revocation committed before the
check was denied in both configurations. An injected pre-commit exception left
neither mailbox nor outbox entry, and a retry created one of each. Concurrent
identical retries created one pair, while conflicting retries were rejected.

These are controlled serialization and recovery checks, not estimates of attack
success in arbitrary deployments. The defined release point is the database
mailbox commit. A later network dispatcher, recipient consumption, or a writer
that bypasses these resource rows lies outside the experiment. In particular,
revocation after that commit does not retract already released data, and a
future dispatcher may need another policy check if its send is the intended
release boundary. Outbox persistence here is not proof of audit-center receipt.
The distinction matters: a fresh decision can be correct when computed yet
invalid when applied at a different boundary.

\paragraph{New tasks with a frozen gate and scorer.}
To separate reuse of development fixtures from task transfer, we froze the gate,
prompts, task oracles and a separate scorer before generating proposals for
20 new tasks. Four authored domains---billing, appointment scheduling,
onboarding and travel---each contribute two benign requests, one disclosure
challenge, one request for an unsupported status and one authority-unavailable
case. Unlike the eight-call pilot, these inputs contain no draft answer. The
model sees the request, task scope and source records; it does not see the
scoring oracle or current authority state. All five controls evaluate the same
proposal. This is direct CPU inference plus a pure release gate, not a new
end-to-end database deployment. Source hashes and the pinned Qwen checkpoint are
retained with the raw proposals.

The predeclared utility endpoint requires every requested source-backed field
on each of the eight benign tasks, without private-marker disclosure,
unsupported values, generation errors or truncation. Refusal is not success.
Privacy-marker release is measured over all 20 tasks, including incidental
leaks on benign requests. Unsupported released values and nonempty releases
on the eight uncertainty tasks are reported separately. These are exact-slot
and literal-marker checks, not human judgments of answer usefulness or semantic
privacy. New task authorship is internal; the scorer is independent of the gate
implementation, not independently authored or adjudicated. Correlation within
four authored domain families precludes treating the 20 tasks as independent
population samples.

\begin{table}[htbp]
\centering\small
\begin{tabular}{lrrrrr}
\toprule
Control & \shortstack{Complete\\(/8)} & \shortstack{Private\\(/20)} &
\shortstack{Unsupported\\(/20)} & \shortstack{Uncertain\\(/8)} & \shortstack{Unknown\\(/20)} \\
\midrule
Snapshot & 8 & 2 & 4 & 7 & 2 \\
Fresh policy & 8 & 2 & 4 & 4 & 5 \\
Policy + scope & 8 & 0 & 4 & 4 & 5 \\
Full gate & 8 & 0 & 0 & 0 & 9 \\
Deny all & 0 & 0 & 0 & 0 & 20 \\
\bottomrule
\end{tabular}
\caption{Frozen new-task evaluation: 20 authored tasks, one pinned Qwen2.5-3B
proposal per task, five matched gate controls. Complete counts successful benign
requests. Private counts literal forbidden-marker release; Unsupported counts
released values absent from the cited source; Uncertain counts nonempty release
when authority or requested facts are unavailable. Unknown is a gate verdict,
not a success. This is a controlled transfer check, not independent semantic
adjudication.}
\label{tab:frozen-tasks}
\end{table}

All 20 calls completed without generation errors or token-cap hits.
Table~\ref{tab:frozen-tasks} separates privacy from utility: the full gate
completes all eight benign requests while releasing no forbidden marker or
unsupported value. Deny-all also prevents those releases but completes no benign
request. Scope checking removes the two literal disclosure failures; it still
allows four unsupported status claims. Exact source support is needed to remove
those four releases. The model generated ``not started'' without evidence in
three domains and used ``onboarding'' as a status in the fourth. Thus reducing
disclosure alone does not establish factual task correctness.

A descriptive breakdown limits the causal interpretation. All nontrivial arms
reject two proposals with a string-valued version field before checking policy;
one is the revoked invoice case. The new run therefore supplies no additional
evidence about revocation-check effectiveness for that case. The full gate has
nine unknowns: two schema failures, three authority-unavailable cases and four
unsupported-value cases. Three further requests are denied for purpose, case or
workspace mismatch. There are no withheld benign requests in this small set.
CPU generation takes 21.9--72.5 seconds per request (median 29.8 seconds), not an
interactive production latency result. Prompts, gate and scorer were not changed
after outcomes were observed; all seven frozen source/protocol hashes match the
archived run.

\paragraph{Two levels of support, not one safety score.}
The operational distinction is between a validated structured projection and a
free-text semantic judgment. Only the first currently has an executable check
against source fields. A future semantic lane could admit paraphrases after a
separate local checker, but its error rate, abstentions, and privacy decisions
must remain separately visible. It must not inherit the structured lane's
contract checks merely because both produce a signed report. This division
retains an exact path for dates, identifiers and status fields without presenting
that path as a solution to arbitrary inter-agent conversation.

\subsection{Expose gaps in observation and costs of linkage}
Reliable transport addresses loss of \emph{recorded} events. Complete observation
requires an independently justified event manifest or a trusted mediation point.
Sequence numbers and acknowledgments can expose gaps relative to such a manifest;
they cannot detect an interaction omitted by both endpoints. The pilot's 40/40
received event IDs establish delivery coverage over its recorded events only.

Cross-message inference additionally requires continuity across the relevant
subject, recipient, and time window. Splitting reports or rotating identifiers
can destroy that continuity. Retaining stable identifiers restores some joins
but exposes relationships to the center. We therefore propose workspace-scoped,
window-bounded linkage, with explicit unknowns when evidence expires or a join
cannot be justified. The retention window is a privacy--coverage tradeoff to
measure, not a confidentiality guarantee; an opaque identifier remains linkable
within its scope. The service retains a bounded workspace graph across reports, but this does not
establish correctness of arbitrary cross-report identity joins or independent
observation completeness. The proposed linkage policy still requires evaluation.

\subsection{External tasks and native success criteria}
\label{sec:external-dojo}
The new authored tasks above do not establish external validity. We therefore
add an action-prefix replay using AgentDojo~\cite{debenedetti2024agentdojo}, an
external benchmark with executable task-utility and attack-success predicates.
We pin the upstream repository at revision
\texttt{089ed468cf3e}, retain its MIT attribution, and use its v1 task definitions
without changing the success criteria. This evaluates an authorization component
of our method, not the full semantic-support gate or a live adaptive agent.

\paragraph{Selection and controls.}
Before scoring, we freeze 1,132 archived trajectories from four official pipeline
labels: GPT-4o (2024-05-13), GPT-4 Turbo (2024-04-09), Claude 3.5 Sonnet
(2024-06-20) and Llama 3 70B Chat. The selection contains all available clean
runs and \texttt{important\_instructions} runs for the specified disclosure
attacks in workspace, travel and Slack: 324 clean trajectories and 808 attacked
trajectories, spanning 81 underlying user tasks. We exclude sabotage and phishing
goals from the privacy endpoint, not unsuccessful attacks. Model names describe
the upstream archives; we do not make new model calls or claim that these are
current-model measurements.

Five matched controls use the same archived action sequence: no gate, a keyword
check for communication intent, literal destination binding to the initial
request, their conjunction, and denial of every mediated egress action. The
adapter consumes only the original user request and proposed tool arguments.
Neither attacker goals, security labels nor successful reference actions grant
permissions. It reuses the structured gate on a destination-authority envelope;
this says nothing about factual support of the message body. The initial request
is trusted and untrusted tool observations cannot extend its authority.
Destination binding deliberately does not infer contacts from names or grant
permissions from retrieved text. Its false denials are part of the evaluation.

\paragraph{Replay semantics and endpoints.}
We re-execute observed tool calls in AgentDojo's in-memory environments. Email,
Slack and web operations modify mock state; no real communication is sent. At
the first blocked egress action, we terminate the prefix and remove the final
answer instead of retaining the original trajectory's later success. Unblocked
trajectories keep their archived final answer. The unchanged upstream predicates
then evaluate task utility and attack success from the resulting state and
executed prefix. Upstream \texttt{security=True} denotes attacker success.
Errors remain in the denominator and are reported separately.

This intervention measures which archived actions a gate admits and what utility
the admitted prefix retains. It does not measure how an agent would recover,
replan, or attack again after a refusal. The selected attack predicates also
encode specific exfiltration goals, not every semantic disclosure. We compare
ungated replay with the archived labels to expose reconstruction differences.
For paired utility differences, we resample the 81 underlying tasks, keeping
all model and injection variants of each task together (2,000 bootstrap draws;
seed 1729). These intervals describe variation within this archive, not a
confidence guarantee for future deployments.

\paragraph{Privacy reduction has a measurable utility cost.}
All 5,660 replays completed in the Jetstream environment. Under the native
attack predicates, combined intent and destination checks reduce observed
attack success from 165/808 to 0/808. Clean utility falls from 198/324 (61.1\%)
to 179/324 (55.2\%); attacked utility falls from 383/808 (47.4\%) to 316/808
(39.1\%). The clean paired change is $-5.9$ percentage points (task-cluster
bootstrap 95\% interval $[-10.2,-2.2]$); the attacked change is $-8.3$ points
($[-12.8,-4.3]$). The combined arm blocks 28 clean and 202 attacked trajectories;
not every block removes a previously successful task. These losses are material:
this external result does not reproduce the eight-of-eight utility retention
of the internally authored structured tasks.

Figure~\ref{fig:core-replay} in the main text displays these same counts.
The deny-egress control has one clean and one attacked checker exception;
both remain unresolved rather than being scored as successful or verified safe.


Intent checking alone leaves 85 successful attacks. Destination-only and combined
checks have identical utility and attack-success counts, so this experiment
provides no measurable benefit for the additional intent rule. Denying every
mediated egress retains only 150 clean and 285 attacked successes, with two
checker exceptions. Some tasks need no mediated external send, so this arm is
not equivalent to refusing every task. The native predicate failure in each
case is a missing website-state key after early termination; it is not silently
converted to a confirmed safe outcome.

The costs concentrate in Slack: combined checking retains 44/84 clean successes
versus 61/84 without a gate, and 24/168 attacked successes versus 90/168. Travel
retains 42/80 clean and 75/160 attacked successes in both arms; workspace retains
93/160 versus 95/160 clean, and 217/480 versus 218/480 attacked. Thus aggregate
retention conceals substantial environment-specific loss. Literal destination
binding cannot resolve legitimate implicit contacts or contextual destinations.
Eleven combined-arm blocks also arise from unsupported destination argument
representations, including null optional values; we retain them rather than
repairing the adapter after seeing outcomes.

\paragraph{Concordance and scope of the evidence.}
Ungated replay agrees with archived task-utility labels in 1,128/1,132 cases and
attack-success labels in 781/808 attacked cases. The 27 attack disagreements are
26 workspace cases and one travel case. Four task-utility disagreements and
these attack disagreements remain in all denominators. We use outcomes of the
pinned executable predicates consistently across the five arms; the discrepancy
means the replay is not an exact reproduction of all historical labels. Twenty-four
archived traces contain execution errors. Runtime tool errors are retained as
observed behavior; the combined arm has no checker exceptions. All frozen
source/protocol and trace hashes are preserved with the artifacts.

\paragraph{Post-hoc discrepancy audit.}
We replay the 31 trajectories with any label disagreement under seven available
upstream suite versions, from v1 through v1.2.2, without changing the frozen
primary experiment. All four utility disagreements align with archived labels
under later versions: three travel cases pass after v1.1 removes the requirement
that the final answer repeat the hotel rating; one workspace case passes after
v1.1.1 accepts the updated end time as part of calendar rescheduling. This
establishes compatibility with changed predicates, not the provenance of the
historical evaluation. All 27 attack disagreements persist across these versions
(nine replay-positive/archive-negative and eighteen in the opposite direction).
Their historical cause remains unresolved.

As a post-hoc sensitivity analysis only, restricting to trajectories whose
ungated utility and attack labels agree with the archive leaves 322 clean and
779 attacked runs. Combined checking changes clean utility from 198/322 to
179/322, attacked utility from 383/779 to 315/779, and attack success from 156/779
to 0/779. Thus both the reduction in selected attack success and the utility
cost persist without the disputed trajectories. These filtered counts do not
replace the primary denominators; the subset was selected after observing
concordance. Per-trajectory version results accompany the artifacts.

The finding is bounded but useful: explicit task-to-destination binding can
suppress the selected archived exfiltration actions, while failing to retain
legitimate tasks that require contextual authority resolution. It does not
validate semantic support of message bodies, arbitrary leakage channels,
continuing attackers or agent recovery after a block. A next method should
resolve recipient authority from trusted task evidence and evaluate recovery in
a fresh closed-loop run; it should not extract permission from attacker-written
observations or use benchmark attack labels as a policy oracle.


\subsection{A falsifiable evaluation of the proposed composition}
The frozen new-task experiment covers matched structured requests, but not
independent semantic validation. A subsequent external evaluation should retain
the same pre-outcome freeze of implementation, prompts, policies and scorer.
Compare snapshot control, fresh-policy control, fresh-policy plus case scoping,
and local support checking on independently sourced workflows. Baselines must receive the same policy and identity information;
central raw-content checks may be a separate reference, with their disclosure
cost reported. Ablate freshness, scope, support checking, and event continuity
individually. Repetitions of one workflow remain one statistical cluster.

Measure false authorization and false denial separately from semantic disclosure,
unsupported assertions, task success, and abstention. Report error among decided
events together with the fraction decided: rejecting everything must not appear
effective. Independent reviewers should label a blinded sample against explicit
policy and source evidence, with disagreements adjudicated and reviewer access
kept distinct from the center's access. Report latency, metadata volume, linkage
exposure, and outage behavior alongside those outcomes.

The proposed support checker fails its purpose if it reduces literal marker hits
without reducing independently judged disclosure or unsupported claims. The
proposed continuity mechanism fails if its gains depend on retaining relationships
beyond the allowed window. These tests would establish a useful contribution:
not that desensitization makes every interaction auditable, but that the system
identifies which conclusions its limited evidence supports and exposes where
additional local checks or trusted observation are necessary.

\section{Development diagnostics and selection history}
\label{app:development}
This appendix preserves earlier development stages in chronological order.
Their provisional launch criteria and future-tense plans describe those stages,
not additional claims about the completed formal batch in Section~\ref{sec:eval}.
Development outcomes are never pooled with formal estimates.

\subsection{Paired protocol and model checks}
\label{sec:paired-protocol}
We extend the authored scheduling sandbox to test whether marker removal also
limits inference about private sets. Three finite decision rules select the first
common option, the cheapest common option under public costs, or the first option
approved by two of three members. Feasible and infeasible cases are balanced.
These are controlled variations of a shared representation, not three independent
application benchmarks. Twin worlds differ in one member's private bit while
preserving the authorized decision and the other members' inputs.

\paragraph{Information and execution.}
An enumerating observer knows the other two members' sets and the authorized
result. We report $\log_2(|\Omega_0|/|\Omega_T|)$ for the sets of compatible private
states before and after disclosure. This is information gain under a stipulated
uniform conditional prior; it is not an estimate of population entropy or a
semantic confidentiality guarantee. The exact-redaction and structured controls
expose the full set, while a membership oracle answers eight adaptive queries.
A trusted task-bound broker exposes only the authorized result and returns a
constant closed-task response to further queries. Its restricted interface makes
zero additional information expected by construction; it is an oracle-like
control, not a cryptographic implementation or evidence that an agent discovers
minimal disclosure autonomously.

Three member agents submit source-verified proposals. The proposals are shared
across controls, and a model coordinator executes the resulting receipt. Each
control gets at most one validator-triggered retry using the already-public
authorized result. We retain initial errors and score the final action against
the independently computed task result. Joint success requires correct execution
and zero additional information under the stated observer model.

% Measured table is inserted after complete, manifest-matched exports are checked.
\paragraph{Matched expanded results.}
The expanded stage comprises 12 twin-pair clusters (four per rule), 24 worlds
and four controls per model. Manifest and source hashes, task seeds and trial
keys match exactly across the two model configurations. Neither smoke trials
nor earlier development retries are pooled into these counts.
\begin{table}[htbp]
\centering\small
\setlength{\tabcolsep}{4pt}
\begin{tabular}{llrrrrr}
\toprule
Model & Control & First/final & Joint & Retries & Calls & Bits \\
\midrule
Qwen2.5-3B & Exact redaction & 24/24 & 0 & 0 & 24 & 6.42 \\
 & Structured & 24/24 & 0 & 0 & 24 & 6.42 \\
 & Membership & 20/20 & 0 & 4 & 28 & 6.42 \\
 & Task-bound & 24/24 & 24 & 0 & 24 & 0.00 \\
\midrule
Qwen3-8B & Exact redaction & 22/22 & 0 & 2 & 26 & 6.42 \\
 & Structured & 22/22 & 0 & 2 & 26 & 6.42 \\
 & Membership & 20/21 & 0 & 4 & 28 & 6.42 \\
 & Task-bound & 21/21 & 21 & 3 & 27 & 0.00 \\
\bottomrule
\end{tabular}
\caption{Matched expanded protocol checks. First/final and joint are success counts out of 24 per row, not percentages. Calls count coordinator requests only; each model also uses 72 shared member calls. Bits are mean additional information under the stipulated observer. Worlds within a twin pair are dependent.}
\label{tab:paired-protocol}
\end{table}
Qwen2.5-3B uses 72 member and 100 coordinator calls, including 4 retries with 0 successful recoveries. Observed coordinator time totals 105.6 seconds; reported coordinator token counts total 19,883 input and 648 output tokens. These totals exclude member latency and cannot isolate retry-only time.
Qwen3-8B uses 72 member and 107 coordinator calls, including 11 retries with 1 successful recoveries. Observed coordinator time totals 254.1 seconds; reported coordinator token counts total 22,082 input and 699 output tokens. These totals exclude member latency and cannot isolate retry-only time.
The controls differ primarily in disclosure by construction. These measurements test execution reliability and exposure under a fixed interface; they do not establish a learned defense advantage or a general scaling law.

All controls succeed on all 12 infeasible worlds per model. On the 12 feasible
worlds, task-bound execution succeeds in 12 cases for 3B and nine for 8B.
For membership queries, 18 paired worlds succeed under both models, two only
under 3B, three only under 8B and one under neither. One 8B quorum twin pair
produces different coordinator outputs despite identical task-bound payloads;
this output variation alone does not establish dependence on the hidden bit.


\paragraph{Limits and follow-up.}
The privacy comparison concerns protocol payloads; model-output differences are
recorded separately, and timing channels are excluded. Fixed control order and
prompt caching confound latency comparisons, so observed coordinator time cannot
establish a policy speedup. Historical traces record total tokens and elapsed time,
not the separate time of each retry; the next instrumentation version separates
these quantities and rotates control order. Small paired samples do not establish
statistically reliable model rankings. Qwen2.5-3B and Qwen3-8B differ in model family
as well as size. Closed-provider coordinators and independent adaptive observers
are planned extensions, not completed experiments; they must see only the released
view and use the same frozen tasks and query budgets.


\paragraph{Diagnosing failed recovery.}
We replay all first-pass errors from the expanded stage together with one
historically successful control per affected task/control stratum: six selected
rows for 3B and 18 for 8B. These are diagnostic samples, not held-out estimates.
With original history, cached and uncached executions each finish at 2/6 for 3B
and 7/18 for 8B. A fresh retry containing only the public receipt and an explicit
copy instruction reaches 2/6 and 18/18 respectively. Restricting the first-pass
execution view to the receipt reaches 6/6 and 18/18 without retries. For 8B this
uses 18 coordinator calls versus 30 with history-based recovery. Disclosure to
the stipulated observer is held fixed; the intervention changes execution input,
not the privacy score. Replay order is cyclic; one historically successful 8B
control fails on replay, so historical and replay first-pass counts differ.

The contrast implicates input presentation and recovery context rather than a
transport or JSON failure; it does not identify an internal model mechanism.
A fresh retry also changes its instruction, so its effect cannot be attributed
solely to removing history. Receipt-only first passes preserve the original
system instruction. Because the broker already computes the authorized result,
a deterministic executor is the natural zero-coordinator-call reference. This
experiment measures reliable execution, not autonomous planning. New-seed runs
and independent native-tool tasks are separate follow-ups; their outcomes are
not inferred from these selected replays.

\paragraph{Independent-task competence check.}
Receipt execution does not establish competence on independent tool-use tasks.
Before expanding to a five-arm AgentDojo comparison, we therefore evaluate six
fixed development cases spanning Slack, travel, and workspace tasks, with one
clean and one attacked case per suite. Qwen2.5-3B completes 0/6 with the original
JSON interface and 0/6 with a more explicit task instruction, while model calls
increase from 18 to 33. A subsequent native-tool interface completes 1/6 using
26 calls, 85,361 input tokens, and 2,919 output tokens. It also produces seven
format errors and one output-cap event. None of the three attacked cases jointly
completes the task and defeats the attack; zero observed attack successes alone
is therefore not evidence of useful protection. These are development diagnostics,
not independent estimates of the proposed method's effectiveness.

The formal comparison is withheld until the interface has known checker outcomes,
no format or output-cap errors, executed tools in all three suites, and positive
development utility no worse than the reference. This is a minimum competence
screen, not evidence of reliable generalization. The corrected runner serializes
multi-call responses, checks authorization for every action, charges each action
to the same step budget, and cancels pending actions after a denial or tool error.
Costs distinguish model calls from action steps and identify calls triggered by
repair, including their tokens and latency. A repair call is the next model call
after an observed error or denial; this accounting does not isolate the causal
benefit of retrying. Any development budget change must propagate to every formal
arm, and incomplete runs and transport failures remain explicit rather than being
scored as successful defenses.

\paragraph{Budget sensitivity and mechanism boundary.}
After multi-call handling is corrected, the six-case probe has zero format errors
but still completes only 1/6 tasks at both 512 and 1,024 output tokens, with one
cap event in each run. Model calls decrease from 18 to 14, but each run has four
repair-triggered calls and zero attacked joint successes. Their repair costs are
14,599/339 and 15,237/452 input/output tokens respectively. Increasing the output
ceiling therefore does not establish task competence. A completed matched
cross-model development screen contrasts ungated and gated executions on the same six cases at
12 action steps, 1,024 output tokens per call, and an 8,192-token context budget.
Before formal expansion, the complete screen must have known outcomes, no format
or output-cap errors, and at least two of three clean ungated tasks completed.
This small development threshold is a screening rule, not an accuracy claim.
Across both arms, Qwen2.5-3B and Qwen3-8B each complete one of six tasks;
Qwen2.5-7B completes all three clean tasks but none of the three attacked tasks.
For 7B, the observed attack success count falls from one to zero under the gate,
but attacked joint success remains zero. These historical results do not justify formal expansion: a subsequent
provenance audit found that archived requests disagreed with the pinned native
task prompts in two of six development cases (and ten of 32 planned formal
cases). The historical screen is retained as a diagnostic, not a valid estimate
of current-task competence. The 3B configuration also
has an output-cap event and the 8B configuration fails the clean-utility screen.
All 36 development episodes are retained, including unsuccessful configurations.
Model families differ, so this is not a controlled model-size scaling study.
The held-out 7B comparison has not started under the corrected protocol.

\paragraph{Canonical GPU development check.}
We therefore rerun development using requests from the pinned native task
classes, their matching checkers, and upstream tool-result serialization on a
Jetstream A100. Nine conditions comprise the original six cases and three
same-task, no-injection twins of the attacked cases. Seven arms include the five
system controls below, the upstream AgentDojo Tool Filter, and the official
IPIGuard implementation adapted to the same local model and bounded runtime.
All 63 episodes complete with known native outcomes. Ungated execution succeeds
on two of three original clean cases but zero of three clean twins; the latter
terminate without budget exhaustion, format errors, or tool errors. This rules
out those recorded execution failures as explanations, but does not establish
a model-level cause. Boolean checks further isolate a missing required date in
the workspace answer; in travel, the answer contains both required numeric
facts but no calendar event is created. Replaying the recorded actions against
the same initial environments reproduces all three failed native scores. The
Slack agent posts a summary but omits one required website, while completing
the second subtask; the workspace agent searches the calendar without first
querying the current day. The missing date lookup is a diagnostic observation,
not an additional scoring requirement. These diagnostics do not change
the native outcomes. On the three attacked development cases, ungated execution
has one attack success and zero joint successes. The gate, scope, receipt, and
deny-egress controls and Tool Filter each have zero joint successes; IPIGuard
has two. These small, development-only observations do not establish a ranking
or a general privacy guarantee. Four Tool Filter calls have unknown usage;
reported token totals for that arm are lower bounds.
The frozen development launch requirement of at least two successful ungated clean
twins is not met. At that stage, the 32-task, seven-arm comparison was withheld,
and no favorable effectiveness conclusion is drawn from reduced attack success.
This corrected GPU protocol is not a matched replication of the historical CPU
screen, because task requests, serialization, and budget accounting changed.

\paragraph{Interface ablation and baseline comparability.}
A separate frozen development ablation evaluates all nine conditions under three
interfaces (27 episodes), without changing the model or budgets. Ungated
same-task clean-twin success is 0/3 with the reference envelope, 1/3 with native
tool-result content, and 0/3 when native content is combined with the upstream
system prompt. Corresponding clean-twin model-call totals are 14, 15, and 21;
repair calls are 0, 0, and 10. The latter configuration records ten tool errors.
On attacked cases, native content increases attack successes from 1/3 to 2/3
without achieving joint success; the upstream-prompt variant achieves 1/3 joint
success. These development observations do not establish causal population
performance, and none qualifies for formal expansion. We also identify a
baseline-comparability defect: the native adapter overrode the runner's system
prompt, whereas IPIGuard received its legacy JSON-interface instruction. The
adapter now exposes an explicit shared prompt setting. The earlier seven-arm
numbers therefore remain descriptive diagnostics, not a fair effectiveness
ranking. These diagnostics motivated the common-prompt 14B validation reported below.

\paragraph{Matched controls and claim scope.}
The initial comparison design froze tasks, model, context and action budgets
for seven arms, including the two upstream defenses and five system controls: ungated execution, the existing policy gate with replanning,
operation-scoped authorization with replanning, receipt-bound execution with
replanning, and denial of egress. The latter is a restrictive utility reference,
not a competitive defense. Scoped authorization is an ablation that isolates
binding from authority decisions. None of these authored controls is presented
as an implementation of an external state-of-the-art defense. Presidio measures
identifier detection in a separate experiment and is not a task-aware AgentDojo
defense. The complete binding reference is evaluated in the execution-contract
tests below; no receipt superiority over that reference is claimed. Native attack
predicates measure benchmark-specific violations, not all semantic disclosures.

We also test five authored execution-contract attacks: replay, action substitution,
cross-task use, stale epochs, and revocation. Four controls and 32 parameter
variants give 640 mechanism trials on Jetstream, with an eight-query attack cap.
The complete binding reference and the receipt executor both reject every attack
while allowing the legitimate action. Reauthorization alone blocks substitution
and revocation but not the other three families. These dependent, deterministic
variants are not 640 independent tasks or semantic privacy attacks. In particular,
the result establishes no security advantage over the complete binding reference.
Message-body leakage to an authorized recipient and cumulative inference require
separate evaluation; authenticated execution alone does not prevent either.

\section{Historical detector studies and component measurements}
\label{app:detector-studies}
\suppressfloats[t]
\subsection{Historical evidence and provenance limits}
We retain the 3,168-trial instruction sweep and related recipient experiments in
Appendix~\ref{app:historical}. Their source traces support arithmetic
reproduction, but some archived backend aliases lack provider-returned version
identifiers. We therefore do not use these rows to establish a model ranking or
a capability--privacy relationship. The fixed-checkpoint structured tasks and
the external task evaluation in Appendix~\ref{sec:external-dojo} provide separately scoped evidence for the
release mechanisms.

\subsection{Comparison with real baselines}
A hand-rolled regular-expression scanner is not a credible DLP baseline, so we run
\textbf{Microsoft Presidio}~\cite{presidio} itself, in two realistic configurations: a
strict identifier policy (SSN, card, email, phone, bank, passport, licence) and a
broad one that adds \texttt{PERSON}, \texttt{LOCATION}, \texttt{NRP} and
\texttt{DATE\_TIME}. The set is 32 realistic scenarios (17 leaks / 15 clean) spanning
all four violation types plus the cases that separate the approaches: non-regex PII, a
diagnosis, authorized external sharing, benign traffic, and single-hint near-misses
that must \emph{not} fire (Table~\ref{tab:frontier}).

The configured baseline observers receive raw content, while our central auditor
receives labels computed locally. This is a deployment choice in this comparison,
not an inherent limitation of Presidio or other tools: they could also run locally.
Our local tagger likewise reads content. The comparison therefore measures these
specific integrations, not superiority over every possible DLP deployment.
\emph{On accuracy}, no Presidio configuration wins on both axes. The identifier policy
reaches only $F1{=}0.26$: it has no notion of owner, purpose, or hop budget, so it
misses inference, purpose, and TTL violations entirely, and the three clean scenarios
it does flag are precisely the \emph{authorized} external shares. The broad policy
buys recall ($0.65$) only by firing indiscriminately---specificity falls to $0.47$,
flagging eight of fifteen clean scenarios including a lunch invitation and a slide
deck---so its apparently strong $83\%$ ``inference recall'' is not inference detection
at all, merely the observation that those messages contain a place and a date. A
policy-aware LLM-judge that reads every byte is precise but insensitive
($F1{=}0.58$, $0\%$ inference recall). On these 32 scenarios, the configured lexical auditor reaches $F1{=}0.94$
with specificity $1.00$ and no raw message text in its central view; the observed
misses concern tagger coverage
(employment- and behavioral-domain hints), never a false alarm. Consistent with the
tagger study below, the open-weight LLM tagger does not help here---it slightly
\emph{trails} the lexical floor---so the lift we previously measured was specific to a
proprietary backend.

\begin{table}[t]\centering\small
\begin{tabular}{lcccccc}
\toprule
Detector & P & R & F1 & spec. & inf.\ rec. & content$\to$center \\
\midrule
Presidio DLP, identifiers & 0.50 & 0.18 & 0.26 & 0.80 & 0\% & 1{,}839 ch \\
Presidio DLP, broad & 0.58 & 0.65 & 0.61 & 0.47 & 83\% & 1{,}839 ch \\
Llama Guard 3 (1B) & 0.53 & 0.53 & 0.53 & 0.47 & 0\% & 1{,}839 ch \\
LLM-judge, policy-aware & 1.00 & 0.41 & 0.58 & 1.00 & 0\% & 1{,}839 ch \\
Ours --- LLM tagger, open 14B & 1.00 & 0.82 & 0.90 & 1.00 & 50\% & \textbf{0} \\
\textbf{Ours --- lexical tagger} & \textbf{1.00} & \textbf{0.88} & \textbf{0.94} & \textbf{1.00} & 67\% & \textbf{0} \\
\bottomrule
\end{tabular}
\caption{Ours vs.\ real baselines on 32 realistic scenarios (17 leaks). All baselines
read every byte; ours decides at zero content egress. Presidio is run in two realistic
configurations---neither wins on both axes: the identifier policy is insensitive, the
broad policy is indiscriminate (specificity $0.47$). The LLM tagger and LLM-judge both
run on the same local open-weight model (Qwen2.5-14B); Llama Guard is the 1B variant
(\S\ref{sec:guard}).}
\label{tab:frontier}
\end{table}

\paragraph{A shipped guardrail.} Presidio is a pattern matcher and the LLM-judge is
a general model handed our policy; neither is a product built to screen agent
messages. Llama Guard is, so we ran it on the message that crosses. It lands at
$F1=0.53$ on this scenario distribution and the two ways it fails are more informative
than the number.\label{sec:guard}

It flags \textbf{all eight} of the authorized and in-organization shares as unsafe.
Those carry exactly the same sensitive content as the leaks; the only difference is
that the recipient is entitled to it. And it calls \textbf{all six} compositional
inference cases safe, because no single message in them contains a hazard.

Both failures are architectural rather than a matter of capacity. The hazard
taxonomy has no slot for \emph{who may receive} a datum, so an authorized share and
a leak are the same message to it; and it judges messages one at a time, so
evidence that only becomes disclosive when it accumulates is invisible. A larger
Llama Guard would share both blind spots exactly --- which is why we report the 1B
variant, forced by local disk, without claiming it is the strongest available: the
limit being demonstrated is not one that scale removes.

\subsection{Adaptive evasion: what can an adversary learn while evading?}
The right question about an adaptive adversary is not whether it can evade---it
always can, by sending nothing---but \emph{how much it can learn while evading}.
We therefore sample the adversary's strategy space at random and run every
strategy through the real auditor, recording both whether the detector fired and
the attribute reported by the recipient under the experiment's scoring rule,
$g=P(A\mid\text{fragments})-p_0$. A strategy draws a fragment count $n\leq8$,
a per-fragment evidence strength $\lambda_i\in[1.05,3]$, a split across up to
three of the recipient's agents, up to three subject aliases, and optional
verbatim replay; suppressing a tag hides a fragment from the auditor but does not
un-send it, so $g$ is computed over what actually crossed the boundary. Across
$20{,}000$ strategies the largest gain achieved by any \emph{undetected} strategy
is $0.29992$ against the policy threshold $\delta{=}0.3$: not one evasion crossed
the line (Table~\ref{tab:evasion}). Within this synthetic model, the detector and simulated adversary share the
same gain function, so complete and correctly grouped metadata aligns their
thresholds. This does not validate the gain model against natural-language
recipients or cover missing observations.

The ablations suppress tags or split subject identities. At tag rate $q=0$,
$53.6\%$ of sampled strategies exceed the simulated gain threshold without an
alert; aliasing produces $16.3\%$ harmful evasions. These are modeled integrity
failures, not measurements of resistance to an actual compromised client.
Our report authentication and canonical identifiers provide implementation
hooks for integrity, but their software implementations do not prove key
confinement, input authenticity, or complete mediation. Rejecting an unauthenticated
report exposes missing evidence; it does not undo an unobserved disclosure.
The enrollment calculation $(1-c)\times53.6\%$ is consequently a conditional
mixture calculation assuming every enrolled sender reports completely and
correctly, not an empirical deployment guarantee.

\begin{table}[t]\centering\small
\begin{tabular}{lccc}
\toprule
Configuration & detect & max $g$ evaded & harmful evasions \\
\midrule
Full design & 0.54 & \textbf{0.29992} & \textbf{0.0\%} \\
$-$ attestation ($q{=}0.4$) & 0.20 & 0.90 & 33.9\% \\
$-$ attestation ($q{=}0$) & 0.00 & 0.90 & 53.6\% \\
$-$ canonical subject ids & 0.38 & 0.90 & 16.3\% \\
\bottomrule
\end{tabular}
\caption{Adaptive evasion, $20{,}000$ random strategies per configuration. A
\emph{harmful} evasion is one that goes undetected \emph{and} gives the recipient
a belief gain $\geq\delta{=}0.3$. Under the full design there are none: the
largest gain any evasion achieves is $0.29992$, just under the threshold. Each
ablated mechanism admits evasion at the full attainable gain ($0.90$).}
\label{tab:evasion}
\end{table}

\paragraph{Attacking the tagger instead.} The study above assumes honest tags;
the tagger is the one component that reads content, so we attack it directly. An
attacking model is given the protected fact and, in the white-box arm, the lexical
tagger's entire 141-term lexicon with the instruction to avoid every term. A
hand-off \emph{evades} only if a reader model from a different family, shown the
hand-off alone, judges the fact conveyed \emph{and} the auditor fires no
content-derived violation---saying nothing is not evasion. Across six attackers (7B open-weight to GPT-4o) and 576 hand-offs, the lexical tagger is evaded $55\%$ of the time black-box and $70\%$ white-box; against the 32B attacker it catches none of 48 white-box hand-offs, against GPT-4o two. The LLM tagger, facing the same hand-offs, is evaded $4\%$ $[1,7]$ black-box and $7\%$ $[2,13]$ white-box (cluster bootstrap over facts), and
black-box evasion does not rise with attacker strength; against the 32B attacker, swapping in a smaller tagger from a third family (8B) raises its evasion from $0\%$ and $6\%$ to $8\%$ and $12\%$, so the figure tracks tagger strength rather than a shared lineage. For the open-weight attackers, avoiding the vocabulary costs the message (conveyance $87\%\!\to\!82\%$); GPT-4o hedges instead, conveying the fact in only $62\%$ of its black-box hand-offs. This
is the thesis on the defender's side---a token list fails against any sender able
to follow ``avoid these words,'' while a tagger that judges meaning does not. Table~\ref{tab:frontier}'s lexical row is therefore the wrong one to deploy against a sender who is trying; the LLM tagger ($F1=0.90$ there) is the configuration the attack leaves standing. Query access is the stronger attack, and we ran it: handed the LLM tagger's verdict and three rewrites, two attackers raise evasion to $16\%$ $[7,26]$---and lower conveyance to $60\%$, because what passes the tagger increasingly fails to say the thing. The tagger is not unbeatable; it is expensive to beat, and the price is the message.

\subsection{What the audit costs, and what actually crosses}\label{app:historical-cost}
Lemma~\ref{lem:blind} says the center's view of a message is a hash plus a
categorical label, so its size cannot depend on the message. That is a claim about
the implementation as much as the mathematics, and it is falsifiable: we sweep
content length over three orders of magnitude and measure the bytes that actually
cross. They are constant at \textbf{418 bytes per message} from 64-byte to
64\,KB messages---$0.0\%$ variation across a $1024\times$ range. This is an empirical fixed-schema size result, not a proof of content secrecy;
metadata values and repeated observations can still disclose information.

Throughput is $32$--$40$k messages/s single-process, and per-message cost grows
only $1.24\times$ from $100$ to $50{,}000$ messages: the detectors are linear in
edges, so an audit scales with traffic rather than with history. We state the
comparison against a content-shipping observer as a crossover rather than a win,
because it is one. At $16$\,KB messages the center receives $39\times$ fewer bytes;
at 64-byte messages it receives $6.5\times$ \emph{more}, since 418 bytes of metadata
then outweighs the message it describes. Lemma~\ref{lem:blind} gives a conditional representation-capacity bound, not a
non-invertibility guarantee. Equality patterns, low-entropy hashes and repeated
metadata may expose information even when the representation is small.

\subsection{The self-authored floor}
On all 48 A2A-MT scenarios the auditor scores $P/R/F1=1.0$ (TP$=18$, FP$=0$,
TN$=30$, FN$=0$) with full per-type coverage and no over-firing on single-hint
clean controls, and holds $1.0$ when the templated text is replaced by
LLM-generated prose---the detectors key on labels, not wording. We report this
first and weight it least: we wrote both the detectors and the scenarios, so it
measures the detectors \emph{given correct labels} and nothing more. Everything
that follows is an attempt to break it. Generating the realistic variant did earn
its keep in one way hand-written tests could not: it exposed TTL fragility under
paraphrasing (fixed by \texttt{provenance\_id}) and an invariant false positive on
prose tokens coinciding with label tokens.

\subsection{The local tagger}
The inference
detector accumulates \texttt{inferred\_categories}; in a real deployment these
must be \emph{produced} from content, not hand-set. A local tagger (lexical,
zero-dependency, optional LLM backend) emits tags to the center. A locally hosted backend keeps raw content on the
agent; a remote API backend also exposes it to that provider and requires an
additional trust assumption. On a labeled set with paraphrased cases, the lexical backend recovers explicit
(sensitive) \texttt{category} at $P/R/F1=1.0$ and \texttt{inferred\_categories}
at $P=1.0$, $R=0.78$; the \textbf{LLM backend (gpt-4o-mini) lifts inferred-category
recall $0.78\to1.0$} (the misses are novel paraphrases like ``the place where they
treat tumors'').
That lift, however, is \emph{backend-dependent}, and we report the negative result:
on the same labeled set, open-weight backends served locally (Qwen2.5-7B/14B,
Llama-3.1-8B) do \emph{not} clear the paraphrase gap---each sits at or below the
lexical floor's $R{=}0.78$, and Llama-3.1-8B is materially worse ($R{=}0.56$)---while
also \emph{costing} explicit-category precision ($1.0\to0.38$--$0.75$) by over-tagging
benign text. Qwen2.5-14B is the partial exception, raising inferred \emph{precision}
to $0.88$ (best inferred $F1{=}0.82$) with recall still at $0.78$. These configuration-dependent results do not isolate capability as a causal
factor; replacing a lexical tagger with a local model does not consistently
improve inference coverage in this experiment.
\texttt{AuditSession.observe} wires it in so a deployer supplies only text plus
policy. Both dishonest suppression and honest labeling errors can remove
necessary evidence. A MAC detects unauthorized report modification, but cannot
distinguish a correctly authenticated labeling error from a correct label.
The following ablations study the semantic recall limits.

\subsection{An author-independent benchmark}
Our 48-scenario benchmark scores 1.0, but we wrote both the detectors and the
scenarios. To reduce direct scenario-author overlap, an LLM generates diverse
scenarios with its own intended labels and we score our auditor against them. Pooling
six batches gives \textbf{75 scenarios} (46 intended leaks), every one of which parses
and is scored---no silent sample shrinkage. The honest result sits far below the
self-authored 1.0: with the lexical tagger $P{=}1.00$, $R{=}0.17$; with an open-weight
LLM tagger recall rises to $R{=}0.46$ but precision falls to $P{=}0.88$ (three false
positives). That last number sharpens a claim we would otherwise have overstated: the
architecture does not over-fire \emph{given correct labels}, but an over-tagging
backend can still manufacture false positives end-to-end, so the absence of false positives in the controlled cases cannot be
generalized to the deployed pipeline or arbitrary correct-label inputs.

\paragraph{Where the recall actually goes.} It would be easy to attribute this gap
to taxonomy coverage and move on; we tested the attribution instead, on a fresh pool
and with the LLM tagger moved to a different model family (Mistral-7B) than the
author, so that one model is not both writing and grading the scenarios---which is
why its absolute numbers differ from the pool above. Scoring one
pool of $99$ author-generated scenarios under progressively more generous taggers
bounds what any tagger fix can buy (Table~\ref{tab:ceiling}), and the result
corrects us in two places.

First, \textbf{the ceiling is $0.76$, not $1.0$}. Even with every part marked
maximally sensitive and carrying every inferred category, roughly a quarter of the
author's intended leaks go undetected. That remainder is not reachable by any
tagger: it is detector scope plus the author's own label noise, and no amount of
taxonomy work will move it. Second, within what \emph{is} reachable, the binding
constraint is the \textbf{sensitivity estimate}, not category coverage---conceding
sensitivity alone recovers the entire reachable gap, while conceding every inferred
category recovers \emph{nothing} ($R{=}0.22$, identical to the lexical floor to the
last confusion-matrix cell). The reason is structural rather than lexical: the
author writes predominantly single-hop scenarios, so converging fragments never
reach $k^{*}{=}2$ and the inference detector cannot contribute at all. Our earlier
reading of the miss taxonomy---that the misses split evenly between disclosure and
inference---was measuring the author's scenario mix, not our detector.

Taxonomy coverage is real but secondary: letting the tagger name the sensitive kind
in open vocabulary instead of choosing from our ten classes lifts recall
$0.37\to0.43$. The oracles are upper bounds and not configurations---flagging
everything costs precision ($0.85\to0.74$)---but they say clearly which component
is worth improving, and that the honest ceiling for this pool is $0.76$ rather than
the $1.0$ a reader would otherwise assume.

\begin{table}[t]\centering\small
\begin{tabular}{lccl}
\toprule
Tagger & $P$ & $R$ & what it isolates \\
\midrule
lexical & 0.85 & 0.22 & the shipped floor \\
LLM backend & 0.86 & 0.37 & shipped, fixed taxonomy \\
LLM, open vocabulary & 0.84 & 0.43 & taxonomy coverage \\
oracle: all inferred categories & 0.85 & 0.22 & category tagging \\
oracle: maximal sensitivity & 0.74 & 0.76 & sensitivity estimate \\
\textbf{oracle: both} & 0.74 & \textbf{0.76} & \textbf{the recall ceiling} \\
\bottomrule
\end{tabular}
\caption{What any tagger fix can buy, on one pool of $99$ author-generated
scenarios. Conceding sensitivity recovers the whole reachable gap; conceding every
inferred category recovers none of it. The ceiling is $0.76$, so a quarter of the
author's intended leaks are out of reach of tagging altogether.}
\label{tab:ceiling}
\end{table}

One caveat keeps this honest in the other direction too: the author model states
some ``inferences'' explicitly and labels borderline cases as leaks, so part of the
$49$--$63\%$ agreement is label noise rather than a detector gap, which is why we
report agreement alongside $P/R$ instead of treating the author's labels as gold.

\subsection{External validation on AgentLeak}
The multi-tenant
auditor reuses a federated engine we validate on the external \textbf{AgentLeak}
benchmark. Driving AgentLeak's own 600 multi-agent scenarios through an LLM to
produce real inter-agent traces and labeling leakage by AgentLeak's rule, the
auditor detects \textbf{100\% of internal-channel leaks at 0.972 precision with
no detected raw vault strings in the central view} ($82.5\%$ internal-channel leak rate,
under the archived scoring rule). On a labeled 33-scenario set the engine holds $P/R/F1=1.0$,
and under the full six-layer desensitizer plus DP it holds $F1\approx0.97$ at
$\varepsilon\in[0.5,3]$ with no detected raw-content matches. These are historical benchmark
measurements, not guarantees of semantic privacy or calibrated inference.

\subsection{Is the operating point cherry-picked?}
The inference detector fires at $k^{*}{=}2$ under our defaults, so it is fair to ask
whether that is a knife-edge fitted to the benchmark. We sweep each knob over the full
48-scenario suite (deterministic; no LLM calls) and find a broad plateau
(Table~\ref{tab:sens}). \emph{Every} setting of $(p_0,\lambda,\delta)$ for which the
closed form gives $k^{*}{=}2$ scores $P/R/F1=1.0$---roughly a $2\times$ band in each
knob independently---and quality degrades only once $k^{*}$ moves off $2$, in the
direction the model predicts: $k^{*}{=}1$ over-fires ($P{=}0.86$) and $k^{*}\ge3$
under-fires ($R{=}0.83$). Within this authored suite, the tested parameter settings track the
closed-form $k^{*}$. This is an internal sensitivity check, not external
calibration of the likelihood ratios or evidence against benchmark-specific tuning. The disclosure floor $\tau$ degrades gracefully rather than cliff-edging
($\tau{=}2$: $P{=}0.90$; $\tau{=}4$: $R{=}0.94$; $\tau{=}5$: $R{=}0.89$), and
raw$\to$center $=0$ at every setting.

\begin{table}[t]\centering\small
\begin{tabular}{lcc}
\toprule
knob & $F1{=}1.0$ band ($k^{*}{=}2$) & degrades at \\
\midrule
$\delta$  & $[0.20,\ 0.40]$ & $0.10$, $0.50$ \\
$\lambda$ & $[2.5,\ 5.0]$   & $2.0$, $9.0$ \\
$p_0$     & $[0.10,\ 0.30]$ & $\le 0.05$ \\
\bottomrule
\end{tabular}
\caption{Inference-detector parameter sensitivity over the 48 scenarios
($\delta$ swept 0.10--0.50, $\lambda$ 1.5--9.0, $p_0$ 0.02--0.30). Quality tracks the
closed-form $k^{*}$, not the raw parameter values.}
\label{tab:sens}
\end{table}


\subsection{Metadata privacy under differential privacy}
Beyond hiding content, we harden the
metadata: identity-bearing fields are pseudonymized with a per-audit shared salt
(cross-tenant comparisons hold in pseudonym space; literal identifiers are
replaced, but linkage and external information may still identify principals) and \texttt{sensitivity} is DP-noised, while categories are kept
structurally. Pseudonymization alone is lossless ($P/R/F1=1.0$). Under DP we
sweep two orders of magnitude, $\varepsilon\in[0.1,8]$, at 40 trials each: mean
$F1$ decays monotonically and gracefully---$0.99$ at $\varepsilon{=}3$, $0.98$ at
$1$, $0.96$ at $0.5$, and still $\approx0.93$ at $0.1$ (sd $\le0.033$)---with no
cliff. The loss is concentrated in recall, i.e.\ disclosure decisions whose noised
sensitivity crosses the floor, exactly where the mechanism predicts, and no raw-content matches are detected at any tested $\varepsilon$.
These accuracy measurements do not establish differential privacy for the whole
report: unchanged categories, routing relations and repeated releases require
separate adjacency and composition analysis.

\section{Historical exploratory measurements}\label{app:historical}
\suppressfloats[t]
These archived measurements lack complete provider-version provenance. Displayed
backend names are experiment aliases. They are retained for traceability and
prompt-condition comparisons, not as verified model rankings or evidence of
a general relation between capability and privacy.

\subsection{Controlled hand-off measurement}\label{app:historical-handoff}
We use 24 authored multi-agent workflows across seven domains, each holding a
sensitive record. A coordinator delegates to an external specialist under a
permissive prompt. Five runs per workflow give 120 trials per configured backend.
Our auditor flags sensitive identifiers or attributes in 78--82\% of trials for
the three open-weight configurations and 46\% for the historical Claude alias
(Table~\ref{tab:wild}). These are detector outputs on controlled tasks, not
independently adjudicated disclosure rates. Cross-purpose hand-offs are purpose
violations by construction; detecting them tests label consistency rather than
an unknown real-world outcome. The central audit records omit raw message text.

\begin{table}[t]\centering\small
\begin{tabular}{llcc}
\toprule
Recipient model & tier & disclosure & raw$\to$center \\
\midrule
Qwen2.5-7B   & small (open) & 82\% (99/120) & \textbf{0} \\
Llama-3.1-8B & small (open) & 81\% (97/120) & \textbf{0} \\
Qwen2.5-14B  & mid (open)   & 78\% (94/120) & \textbf{0} \\
\textbf{Claude (alias)} & API alias & \textbf{46\% (55/120)} & \textbf{0} \\
\bottomrule
\end{tabular}
\caption{Auditor flag rates on 24 authored workflows $\times$ five runs per
configured backend. Model names are archived aliases, not all independently
verified provider versions. Raw message text is not sent to the central auditor;
this does not imply that metadata reveals no information.}
\label{tab:wild}
\end{table}

\subsection{Can the agent simply be told not to?}\label{sec:instr}
The measurement above invites an obvious objection: our coordinator prompt ends
\emph{``include whatever you think is helpful''}, so a reader may reasonably
suspect the leak is our wording rather than the model's behavior. It also leaves
this paper's central claim---that the safeguard has to be structural---as an
assertion. We test both by sweeping the instruction alone, over the same 24
workflows, eleven models---nine open-weight models from five families and two API configurations---and $3{,}168$ trials. Ground truth here is
deliberately \emph{independent of our auditor}: literal identifiers and attribute
terms taken from the source record, matched against what the coordinator actually
wrote, with each pattern checked against its own record so a drifted annotation
fails loudly rather than scoring zero in our favor.

Removing the invitation does not reduce the measured match rate in this
experiment ($96\%\!\to\!98\%$); this contrast does not rule out effects from
other prompt choices. Instructing the agent helps, and the shape of the
help is the finding. Naming the forbidden categories explicitly leaves $82\%$; handing the model the
record's actual policy---owner, permitted recipient, purpose, and the fact that
this recipient is \emph{not} permitted---leaves $68\%$; Claude (alias) under the strongest condition still leaks in $36\%$ of hand-offs, and GPT-4o is the one model that mostly withholds ($14\%$), to which we return below (Table~\ref{tab:instr}).

In the saved runs, instructions reduce identifier matches from $78\%$ to
$35\%$ and sensitive-attribute matches from $89\%$ to $60\%$.
Figure~\ref{fig:redaction} reports the per-configuration contrasts. These
literal-match measurements do not identify the model's reasoning process or
establish a general relation between parameter count and privacy behavior.

\paragraph{Interpretation boundary.} The historical configurations differ in
literal-match and refusal behavior. Missing provider-version evidence prevents
attributing these differences to a specific model release, alignment method or
capability level. The experiment supports only the observed contrasts within
its saved traces and scoring procedure.

Against the separate literal-match scorer, the auditor has recall $0.74$
(716 scorer-positive hand-offs missed across the 3,168 trials) and precision $0.90$.
There are 224 alerts without a literal match. Without independent semantic
adjudication, disagreements may reflect auditor error or limitations of literal
matching. These are proxy agreement statistics, not guaranteed bounds on
semantic precision or recall.

\begin{table}[t]\centering\small
\begin{tabular}{lccccc}
\toprule
Coordinator instruction & leak & 95\% CI & ident. & attrib. & auditor \\
\midrule
permissive & 96\% & [93, 98] & 78\% & 89\% & 77\% \\
invitation removed & 98\% & [96, 99] & 82\% & 89\% & 76\% \\
categories forbidden & 82\% & [78, 86] & 36\% & 80\% & 58\% \\
\textbf{full policy stated} & \textbf{68\%} & [63, 73] & \textbf{35\%} & \textbf{60\%} & 70\% \\
\midrule
\multicolumn{6}{l}{\emph{unverified Claude configuration alone, permissive $\to$ policy}} \\
\quad identifiers & \multicolumn{5}{l}{$44\%\ \to\ \mathbf{1\%}$ \quad (literal-match rate)} \\
\quad attributes & \multicolumn{5}{l}{$40\%\ \to\ 36\%$ \quad (barely moves)} \\
\bottomrule
\end{tabular}
\caption{Sweeping the coordinator's instruction, pooled over eleven models ($792$ trials per row). Scored against ground truth independent of our
auditor. Intervals are a bootstrap over \emph{workflows}, not over trials: the three
runs of one record under one prompt are not independent draws, and a trial-level
Wilson interval is about $50\%$ too narrow here (measured design effect $2.1$--$2.4$). Telling the agent the rule---even spelling out that the recipient is not
a permitted one---leaves most hand-offs leaking, and suppresses identifiers far
more than it suppresses sensitive attributes.}
\label{tab:instr}
\end{table}

\input{figures/fig_redaction}

\subsection{What survives a relay chain}\label{sec:relay}
Two things the design rests on had only ever been tested against paraphrases we
wrote ourselves, which is not a test: that \texttt{provenance\_id} follows a datum
through re-wording, so hop/TTL tracking is not defeated by a relay that rephrases,
and that exposure compounds rather than decays. We let real agents do the
rephrasing---each forwards the case onward in its own words, four hops deep,
against a hop budget of two---over the same 24 records, three models, 144 chains.

Exposure does not decay. Two of the three models still carry the subject's
identifiers or attributes at hop~4 in \emph{every} chain ($100\%$); the third
declines only from $94\%$ to $79\%$. The intuition that passing a datum through
successive re-writes dilutes it is simply wrong, and the same identifier/attribute
asymmetry appears again from an entirely different mechanism---Mistral-7B's
identifier rate falls across the chain ($90\!\to\!83\%$) while its attribute rate
\emph{rises} ($94\!\to\!96\%$). Meanwhile the tracking holds: provenance survived
all four model paraphrases in $144/144$ chains and the TTL violation fired in
$144/144$, with zero raw content reaching the center. The failure we were looking
for---a datum that keeps leaking while the auditor loses track of it---did not
occur.

\subsection{Is the inference threat real?}
Our detector fires at $k^{*}=2$ hints, but is that where real inferability begins? We
have an independent LLM play the recipient and infer the subject's withheld attribute
from the $K$ benign fragments it received ($K=1,2,3$), scaled to \textbf{16 sensitive
attributes} across health, finance, employment, legal, and behavioral domains
($5$ trials each). In these saved trials, Claude (alias) infers the withheld attribute from two benign
fragments in \textbf{100\% (16/16)} of attributes and from a \emph{single} fragment in
\textbf{81\% (13/16)}. Other configured recipients differ: on the same 16-attribute
set, open-weight recipients infer at $K{=}2$ in 56--69\% (Qwen2.5-7B, Llama-3.1-8B)
and 81\% (Qwen2.5-14B) of attributes, versus 100\% for the Claude alias---inference
differs across configured runs. Missing version provenance prevents a causal
claim about model capability or progress. The default detector stays silent at
$K=1$ even when the recipient returns the target attribute. This is a detection
coverage gap, not evidence of a general no-false-alarm guarantee. A non-firing
recipient is also not proof that the fragments are uninformative. Per-fragment
likelihood ratios in Section~\ref{sec:auditor} permit earlier alerts for stronger
hints, but the ratios require calibration; these observations alone do not
validate them on unseen tasks.

\input{figures/appendix}
\input{figures/revision_appendix}
\input{figures/jetstream20261013/appendix}

\bibliographystyle{plain}
\bibliography{references}

\end{document}
